\section{Transfer to native Markov edits}\label{sec:native}

\subsection{The structure of an attained cover}

\begin{proposition}[Effect types of elementary edits]\label{prop:effects}
Every elementary edit has exactly one of the effect types in
\Cref{tab:effects}. The types are determined by the named experiments.
\end{proposition}

\begin{table}[ht]
\centering
\small
\begin{tabularx}{\linewidth}{@{}lX@{}}
\toprule
Effect & Change in the experimental record\\
\midrule
Local rewrite & Vary the prepared law at fixed profile within the prescribed tolerance.\\
Gate & Contract support by a cover at fixed projected marked law.\\
Route & Advance the timed license catalogue by a cover at fixed law.\\
Earlier disclosure & Refine the available decision partition at one earlier grade by a cover at fixed law.\\
Terminal packaging & Refine the terminal decision partition by a cover at fixed law.\\
Publication & Advance the audit catalogue by a cover at fixed law.\\
\bottomrule
\end{tabularx}
\caption{The measured distinctions in the elementary edit relation.}
\label{tab:effects}
\end{table}

\begin{proof}
The fixed-profile case is the local rewrite. Now let \(p<q\) be an
attained cover and choose lawful laws at \(p\) and \(q\).
If their supports differ, use the register table from \(q\) at the
rates and spectrum of the law at \(p\), while keeping the old grades,
licenses and publications. The later register satisfies the old coarser
conditional decision tests. Rates and spectrum continue to pass their
old tests. This realizes the intermediate profile obtained by replacing
the old support alone with the new support. It lies strictly above \(p\).
The cover makes it equal to \(q\).

Suppose the support agrees. The later law can realize an advance in
licenses alone by taking the full license set of \(q\) and all other
fields of \(p\). The policy holds for that license set by attainment of
\(q\), even if the policy would deny some intermediate subsets.
An advance in publications alone is realized by the same later law
at the old license set. Thus if either catalogue differs, the
corresponding intermediate profile already equals \(q\). The update
may add several names when the cover skips unattained intermediate
sets.

In the remaining case only the grades differ. Advance the latest
differing grade to its value in \(q\). Earlier grades were coarser than the
old value at that grade, and every later grade already agrees with \(q\).
Nesting is therefore preserved. The later law realizes this profile
since each of its new decision partitions is no finer than that of \(q\).
The cover property again forces the intermediate profile to be \(q\).
The changed grade itself distinguishes the earlier and terminal cases.
All the cases are disjoint, and their tested fields are recoverable
from the experiments.
\end{proof}

\subsection{The metric return and the rate fiber}

Let \(N=|X\times R|\). Sum one prepared row of \(P_m\) to recover
\(\theta_m\). Use \eqref{eq:signal} to recover the response from the same
prepared row by dividing by \(\theta_m\).
If the native sup-norm difference between two laws is \(\eta\),
each of these two numerators changes by at most \(N\eta\), since
the signal weights are signs. A reference response has magnitude at
most one, and both rates are at least \(\delta\). Subtracting the
quotients gives
\begin{equation}\label{eq:inverse-modulus}
 \Norm{\alpha_A-\alpha_B}\le \frac{2N}{\delta}\,\eta .
\end{equation}
Conversely, the polynomial formulas \eqref{eq:native} are uniformly
continuous in the rates, responses and register entries on the
bounded boxes under consideration. One can use the bound on the
difference of three-factor products by the sum of their differences,
since \(\theta_m,Q_m,L_m\) all lie in \([0,1]\).
These observations connect numerical steps in both directions.

For use in the forward construction, the rate laws at a fixed
profile are finely connected. To see this, let \(\theta,\theta'\)
be two feasible rational rate vectors. If they agree, keep them fixed.
Otherwise take the componentwise geometric interpolation
\(\gamma_m(t)=\theta_m^{1-t}(\theta'_m)^t\) and normalize by
\(\sum_m\gamma_m(t)\). Homogeneity preserves every test after
normalization. Every numerator is at least \(\delta\).
For \(0<t<1\) the sum is strictly less than \(1\), by the
weighted arithmetic--geometric mean inequality and the fact that at
least one coordinate changed. All normalized rates are therefore
strictly above the floor for interior times. Use the prime-exponent
construction of \Cref{lem:rational-chart} for the positive endpoint
ratios, then normalize the rational interpolants. This keeps the tests
exact and works sufficiently close to every interior mesh point
because the floor has slack. The endpoints themselves are used
verbatim.

\subsection{Proof of native reachability}

\begin{proof}[Proof of \Cref{thm:native}]
For the forward direction, every elementary edit increases the profile
in \eqref{eq:profile-order} or keeps it fixed. Each later law therefore
passes all initially held spectral tests, and all discrete covers preserve
the spectrum. For any prescribed positive rational spectral mesh
\(\rho\), take a native chain at tolerance
\(\delta\rho/(2N+1)\). By \eqref{eq:inverse-modulus} its projected spectral
steps have size less than \(\rho\). \Cref{thm:spectral} then gives equality
of the signed cores relative to the initial tests.

For the converse assume profile comparability and equality of these
cores. Form the intermediate profile \(p'\) by taking the support
of the target profile and all other fields of the source. It is attained
by the target register at the source rates and spectrum. Between the
source profile and \(p'\), the other fields are fixed. Follow a finite
cover chain in that interval, selecting a register with the attained
support at each step and keeping the source rates and spectrum. These
are legal gates because the register constraints and held rate and spectral
constraints have separate variables.

At \(p'\) the current and target laws both meet the source's tests.
The two register tables have the same positive support. Their affine
interpolation preserves the prescribed margins, reciprocal pairing and
conditional block masses. It is positive exactly on that support, and a
rational mesh realizes its changes. Keeping the final register fixed, move
the rates by the preceding rate argument and move the spectrum by
\Cref{thm:spectral}, both with the source's tests held. The three numerical
chains can use a small enough mesh to meet the given native tolerance
by \eqref{eq:native}. At this stage the target law is attained with profile
\(p'\). Follow a finite cover chain from \(p'\) to the target profile
while keeping that law fixed. Every intermediate attained profile has
the target support and weaker tests. Its license policy holds because
the profile is attained. This makes each step legal and yields the
requested chain for every positive rational native tolerance.
\end{proof}

\begin{proposition}[Readback of edits]\label{prop:readback}
For a fixed instrument, the prepared probabilities and the named
decision, license and publication responses determine the state.
The correspondence with attained query tables preserves and reflects
elementary edits, effect types and finite chains, including chains with
a specified effect omitted. A relabeling that identifies the full
attained profile order and the named queries gives the same conclusion
for the relabeled instrument.
\end{proposition}
\begin{proof}
The positive entries give the support, and the available decisions give each
partition by their common equality relation. Catalogue membership is read
from the remaining queries. The prepared probabilities determine rates and
spectra by the preceding row-sum and signal calculations. Finally
\[
 L_m(g,h;u,v)=\frac{P_m((g,u),(h,v))}{\theta_m Q_m(g,h)}
\]
recovers every register entry. The denominator is strictly positive by
\eqref{eq:guards}. Thus attained states and their attained query tables
correspond in both directions. The fixed-profile comparison, cover order and
law or projected-law preservation in the edit definition can all be tested
on the tables. The same tests define the edits after relabeling when
the relabeling identifies the attained order and the experiment labels,
including the terminal grade. Lifting each step in a finite chain proves
the claim.
\end{proof}
